\chapter{退化圆锥曲线}

\section{二次点与二次线}

回顾圆锥曲线的定义。

\begin{definition}[圆锥曲线]
一个不定的对称可逆的 $(0,2)$ 型张量 $g_{ij}$ 称为\textbf{圆锥曲线张量}。点 $a$ 在圆锥曲线 $g$ 上由
\[
g_{ij}a^ia^j=0
\]
刻画。$g$ 的逆记为 $(g^{-1})^{ij}$，当把 $g$ 取作度规时简写为 $g^{ij}$；直线 $l$ 与 $g$ 相切由
\[
(g^{-1})^{ij}l_il_j=0
\]
刻画。
\end{definition}

圆锥曲线的定义要求 $g_{ij}$ 可逆，但不可逆的张量同样有重要的用处。例如双曲线
$\frac{x^2}{a^2}-\frac{y^2}{b^2}=k$ 在 $k=0$ 时成为两条直线，这两条直线可以看作一种特殊的圆锥曲线，
称为\textbf{退化圆锥曲线}。由点与线的对偶性，“两个点”也属于这一类：它是被压扁成一条线段的椭圆的
两个端点。为了把“两条直线”“两个点”与圆锥曲线放在同一个框架中研究，需要引入比圆锥曲线更一般的
张量。

另一个基本事实是：一条直线与一条圆锥曲线交于两点 $a,b$，或者过一个点作圆锥曲线的两条切线，这
两个点（或两条线）具有完全相同的地位，不能拆开而单独看待。用某个不可逆的二阶张量（而不是用
两个点矢或两个线矢）表示这样的一对对象，是最自然的做法。

\begin{definition}[二次点与二次线]
设 $f^{ij}$ 是不定对称的 $(2,0)$ 型张量，$h_{ij}$ 是不定对称的 $(0,2)$ 型张量。
\begin{enumerate}[label=(\arabic*),leftmargin=2.4em]
\item $f^{ij}$ 称为\textbf{二次点张量}，它所对应的几何元素称为\textbf{二次点}；有时也直接把
$f^{ij}$ 称为二次点。
\item $h_{ij}$ 称为\textbf{二次线张量}，它所对应的几何元素称为\textbf{二次线}。
\item 二次点与二次线统称\textbf{二次型}，或统称\textbf{二次曲线}。
\item 若 $f^{ij}l_il_j=0$，称直线 $l$ \textbf{经过}二次点 $f$。满足此方程的所有 $l$ 称为一个
\textbf{二次线束}；当 $f$ 不可逆时，称为\textbf{退化二次线束}。
\item 若 $h_{ij}a^ia^j=0$，称点 $a$ 在二次线 $h$ 上。满足此方程的所有 $a$ 称为一个
\textbf{二次点列}；当 $h$ 不可逆时，称为\textbf{退化二次点列}。
\end{enumerate}
\end{definition}

有一个上指标的量是点，有两个对称上指标的量是二次点；有一个下指标的量是线，有两个对称下指标
的量是二次线。因此“二次点”“二次线”正是点与线的推广。

\begin{remark}
经过可逆二次点 $f^{ij}$ 的直线，就是与圆锥曲线 $f$ 相切的直线，故“经过二次点”确实是“经过一个
点”的推广。二次型可以可逆，也可以不可逆：圆锥曲线 $g_{ij}$ 是二次线的可逆特例，它自然地生成
二次点 $(g^{-1})^{ij}$；而不可逆的二次线不能生成一个二次点，反之亦然。
\end{remark}

\section{不可逆二次型的几何意义}

按二次型矩阵的秩，也就是号差中 $0$ 的个数，可以把二次型分为三类：没有 $0$、有 $1$ 个 $0$、
有 $2$ 个 $0$。没有 $0$ 意味着可逆，对应圆锥曲线；有 $1$ 个或 $2$ 个 $0$ 的二次型不可逆。
下面证明：号差中有 $1$ 个 $0$ 的二次型对应“两个点”或“两条线”，有 $2$ 个 $0$ 的二次型对应
“单个点”或“单条线”。以二次线为例给出计算，二次点的情形完全对偶。

\begin{proposition}
\label{prop:onezero}
设二次线 $h_{ij}$ 的号差为 $(1,1,0)$，即恰有一个 $0$。
\begin{enumerate}[label=(\arabic*),leftmargin=2.4em]
\item 存在非零的点矢 $s$，使得 $h_{ij}s^j=0$；满足 $h_{ij}x^j=0$ 的所有 $x$ 都与 $s$ 共线，
即 $x^i=\lambda s^i$。并且 $h_{ij}s^j=0$ 不仅意味着 $h_{ij}s^is^j=0$，还意味着对任意矢量 $a$
都有 $h_{ij}s^ia^j=0$。
\item $h$ 上点的全体是过 $s$ 所对应的点 $S$ 的两条直线。
\end{enumerate}
\end{proposition}

\begin{proof}
号差为 $(1,1,0)$ 意味着 $\rank h=2$，故映射 $x\mapsto h_{ij}x^j$ 的核是一维的：存在非零的
$s$ 使 $h_{ij}s^j=0$，并且解空间由 $s$ 张成。由 $h$ 的双线性性，$h_{ij}s^j=0$ 立刻给出
$h_{ij}s^ia^j=0$ 对一切 $a$ 成立。

由 $h_{ij}s^is^j=s^i(h_{ij}s^j)=0$ 知 $S$ 在 $h$ 上。设 $A$ 也在 $h$ 上，即
$h_{ij}a^ia^j=0$。直线 $SA$ 上的点都可写成 $s+\lambda a$，而
\[
h_{ij}(s+\lambda a)^i(s+\lambda a)^j
=h_{ij}s^is^j+2\lambda h_{ij}s^ia^j+\lambda^2h_{ij}a^ia^j=0,
\]
因此整条直线 $SA$ 都含于 $h$。于是 $h$ 上点的全体是过 $S$ 的若干条直线之并。

下面证明这样的直线恰有两条。先取一条不经过 $S$ 的直线 $\ell$，$\ell$ 上的点都可写成
$p+\lambda q$。多项式
\[
P(\lambda)=h_{ij}(p+\lambda q)^i(p+\lambda q)^j
\]
关于 $\lambda$ 至多是二次的。若 $P\equiv0$，则 $\ell$ 上每一点都在 $h$ 上，把 $\ell$ 上任一点
与 $S$ 相连即知过 $S$ 的每条直线都含于 $h$，从而平面上每一点都在 $h$ 上，即
$h_{ij}a^ia^j=0$ 对一切 $a$ 成立，这与 $h\neq0$ 矛盾。因此 $P$ 至多有两个零点，即 $\ell$ 与
$h$ 至多交于两点。而 $\ell$ 上每一点若在 $h$ 上，就对应过 $S$ 的一条含于 $h$ 的直线（这是因为
$S$ 与 $\ell$ 上该点都在 $h$ 上，连线整条含于 $h$）。所以过 $S$ 且含于 $h$ 的直线至多两条。

再证这样的直线恰恰有两条。把 $h$ 对角化，在适当的基下 $h$ 的分量为
$\diag\{1,-1,0\}$。由前两个基矢张成的二维子空间中，$h$ 的限制为 $x^2-y^2$，它在两条直线
$x=\pm y$ 上取零值。这两个方向与核方向（第三个基矢）不共线，故对应的两个点都在 $h$ 上且都不
等于 $S$，从而过 $S$ 的这两条直线都含于 $h$。故过 $S$ 且含于 $h$ 的直线恰有两条。
\end{proof}

\begin{proposition}
\label{prop:twozero}
设二次线 $h_{ij}$ 的号差中有恰好两个 $0$。则存在两个不共线的点矢 $s,t$ 使
$h_{ij}s^j=h_{ij}t^j=0$，并且 $h$ 上点的全体恰为 $s,t$ 所连的直线。
\end{proposition}

\begin{proof}
号差中有两个 $0$ 意味着 $\rank h=1$，映射 $x\mapsto h_{ij}x^j$ 的核是二维的，取其中的一组基
$s,t$，就有 $h_{ij}s^j=h_{ij}t^j=0$，且 $s,t$ 不共线。由 $h_{ij}s^is^j=h_{ij}t^it^j=0$ 知
$s,t$ 都在 $h$ 上；与上面相同的计算表明直线 $st$ 含于 $h$。

设 $p$ 不在直线 $st$ 上，则 $p,s,t$ 构成三维空间的一组基，任一矢量 $a$ 可写成
$a=\lambda p+\mu s+\sigma t$。由 $h_{ij}s^j=h_{ij}t^j=0$ 得
\[
h_{ij}a^ia^j=\lambda^2h_{ij}p^ip^j .
\]
若 $p$ 也在 $h$ 上，则对一切 $a$ 都有 $h_{ij}a^ia^j=0$，即 $h=0$，矛盾。因此不在直线 $st$ 上的
点都不在 $h$ 上，$h$ 上点的全体恰为直线 $st$。
\end{proof}

\begin{remark}[号差中 $0$ 的个数与解空间维数]
容易把两个不同方程的维数弄混：线性方程 $h_{ij}x^j=0$ 的解空间维数等于 $h_{ij}$ 的零特征值
个数，即号差中 $0$ 的个数。因此号差中有 $2$ 个 $0$ 时解空间是二维的，存在两个不共线的解
$s,t$；号差中有 $1$ 个 $0$ 时解空间才是一维的。对后一情形，上面的证明用“二次型在一条直线
上的限制若非恒零则至多有两个零点”这一事实补足了“含于 $h$ 的直线恰有两条”的论证。
\end{remark}

\begin{theorem}[退化二次型的几何意义]
在二次线 $h_{ij}$ 上的点构成
\begin{enumerate}[label=(\arabic*),leftmargin=2.4em]
\item 圆锥曲线，当 $h$ 可逆时；
\item 两条直线，当 $h$ 的号差中有一个 $0$ 时；
\item 一条直线，当 $h$ 的号差中有两个 $0$ 时。
\end{enumerate}
对偶地，经过二次点 $f^{ij}$ 的直线构成
\begin{enumerate}[label=(\arabic*'),leftmargin=2.6em]
\item 圆锥曲线的全部切线，当 $f$ 可逆时；
\item 经过两个定点的全部直线，当 $f$ 的号差中有一个 $0$ 时；
\item 经过一个定点的全部直线，当 $f$ 的号差中有两个 $0$ 时。
\end{enumerate}
\end{theorem}

\section{用点矢与线矢表示退化二次曲线}

前面的讨论把退化二次型与点、直线联系了起来。下面给出代数上的联系：如何用（一个或两个）点矢、
线矢把二次型 $f,h$ 具体表示出来。

\begin{proposition}[两个线矢的对称积]
设 $m,n$ 是不共线的两个线矢，令
\[
h_{ij}=m_in_j+m_jn_i.
\]
则 $h$ 的号差为 $(1,1,0)$，它上面的点集恰为直线 $m$ 与直线 $n$ 之并；并且
$s=m\times n$ 满足 $h_{ij}s^j=0$。
\end{proposition}

\begin{proof}
对任意点 $a$，
\[
h_{ij}a^ia^j=(m_in_j+m_jn_i)a^ia^j=2(m\cdot a)(n\cdot a).
\]
因此 $h_{ij}a^ia^j=0$ 当且仅当 $m\cdot a=0$ 或 $n\cdot a=0$，即点 $a$ 在直线 $m$ 或直线 $n$ 上。
于是 $h$ 上点的全体是两条直线 $m,n$ 之并。又 $m\cdot s=n\cdot s=0$，故
\[
h_{ij}s^j=m_i(n\cdot s)+n_i(m\cdot s)=0 .
\]
由 $h\neq0$ 且其零点集为两条直线，$h$ 的号差为 $(1,1,0)$。
\end{proof}

对偶地，对两个不共线的点矢 $a,b$，有
\[
f^{ij}=a^ib^j+a^jb^i,
\]
它是号差为 $(1,1,0)$ 的二次点，经过它的直线恰为过 $a$ 或过 $b$ 的直线之并，并且
$f^{ij}l_il_j=2(a\cdot l)(b\cdot l)$。

\begin{proposition}[一个线矢的平方]
设 $l$ 是一个非零线矢，令 $h_{ij}=l_il_j$。则 $h$ 的号差中有两个 $0$，它上面的点恰为直线 $l$
上的点。
\end{proposition}

\begin{proof}
对任意点 $a$，
\[
h_{ij}a^ia^j=l_il_ja^ia^j=(l\cdot a)^2,
\]
故 $h_{ij}a^ia^j=0$ 当且仅当 $l\cdot a=0$，即 $a$ 在直线 $l$ 上。$h$ 的秩为 $1$，号差中有两个
$0$。
\end{proof}

对偶地，$f^{ij}=a^ia^j$ 是号差中有两个 $0$ 的二次点，经过它的直线恰为过点 $a$ 的直线。
习惯上把二次线 $l_il_j$ 写成 $ll$，把二次点 $a^ia^j$ 写成 $aa$。

\begin{remark}
点矢与线矢可以用上述特殊的二次点与二次线来表示，但一般的二次点与二次线还包含圆锥曲线等丰富得多
的几何对象。可以在这些对象之间赋予代数运算，以表达多样的几何关系；这些运算的内容多于矢量运算，
矢量之间的点乘与叉乘只是其特例。
\end{remark}

例如，点 $a$ 在直线 $l$ 上等价于
\[
(aa):(ll)=a^ia^jl_il_j=(a^il_i)^2=0,
\]
而点 $a$ 在圆锥曲线 $g$ 上等价于
\[
(aa):g=g_{ij}a^ia^j=0 .
\]
前者是后者的特例。

\section{省略指标的写法}

今后会遇到复杂的表达式，省略指标常常更便于书写。约定如下。

\begin{enumerate}[label=\textbf{(\arabic*)},leftmargin=2.6em]
\item 两个元素直接并列，表示它们的指标直接并列而不缩并。例如 $ga$ 的分量为 $g_{ij}a^k$。
\item 两个元素之间加一个点，表示缩并相邻的一个指标。例如 $g\cdot a$ 的分量为 $g_{ij}a^j$，
$f\cdot g\cdot a$ 的分量为 $f^{ij}g_{jk}a^k$。
\item 两个元素之间加一个冒号，表示缩并相邻的两个指标；这一般只用于一个二次点与一个二次线之间。
例如 $f:g=f^{ij}g_{ij}$ 是一个实数。$h_{ij}$ 与 $g_{ij}$ 之间可以进行两种缩并：
\[
h^{-1}:g=g:h^{-1}=(h^{-1})^{ij}g_{ij},\qquad
g^{-1}:h=(g^{-1})^{ij}h_{ij}.
\]
注意这两个数一般不相等；其中一个为 $0$ 时，另一个未必为 $0$。
\end{enumerate}

有了这套写法，$f^{ij}=a^ib^j+a^jb^i$ 可以写成 $f=ab+ba$。由于这个构造十分重要，单独引入记号
\[
f=(ab),\qquad (ab)^{ij}=a^ib^j+a^jb^i,
\]
称为两个点的\textbf{对称积}。同理，两条线的对称积是
\[
h=(mn),\qquad (mn)_{ij}=m_in_j+m_jn_i .
\]

\section{两个点不能拆开}

二次点 $(ab)$ 中的两个点 $a,b$ 具有完全相同的地位，不能拆开而单独看待。这句话的意思是：虽然
号差为 $(1,1,0)$ 的二次点 $f$ 唯一对应了两个点 $a,b$，却很难把这两个点分别表示出来。

已知 $a,b$ 时，可以用与坐标系无关的张量运算写出它们的对称积 $f=(ab)$。反过来，不可能仅仅用
张量积、缩并的语言由 $f$ 表示出 $a$（或 $b$）。这是因为张量积与缩并都是单值函数。若存在单值
函数 $F$ 使 $a^i=F(f^{ij})$，由于 $a,b$ 在 $f$ 中地位平等，同样应当有 $b^i=F(f^{ij})$，
于是 $a^i=b^i$，矛盾。事实上，要用 $f$ 表示 $a,b$ 必须开根号，而开根号的一个正根与一个负根
恰好分别给出 $a,b$。使用张量语言的初衷是仅用张量运算处理几何问题，对分量的函数开根号破坏了
这一点。这使得在很多情况下，只能整体地处理两个点的对称积，而不能对它们分别处理。

与两个点的情形不同，对于由 $a^{ij}=a^ia^j$ 生成的二次点，可以由 $a^{ij}$ 表示出点 $a$：任取
一个线矢 $l_i$，有
\[
l_ia^{ij}=l_ia^ia^j=(l\cdot a)a^j,
\]
其中系数 $l\cdot a$ 只要不为 $0$ 就不影响 $a$ 所代表的点。也就是说，若已知点 $A$ 对应的二次点，
只需把它与一条不经过 $A$ 的直线缩并掉一个指标，就得到了 $A$ 的点矢。

\section{联立二次点}

\begin{example}
\label{ex:joined}
一条直线 $l$ 与圆锥曲线 $g$ 交于 $A,B$ 两点，求由 $A,B$ 确定的退化二次线束，即求二次点
$f=(ab)$。
\end{example}

\begin{proof}[解]
设直线 $m$ 在所求的退化二次线束上，则 $m$ 经过 $A$ 或 $B$，即 $m$ 与 $l$ 的交点是 $A$ 或 $B$，
也就是这个交点在 $g$ 上。两直线 $m,l$ 的交点满足
\[
(m\times l)^2=0,
\]
而
\[
(m\times l)^2=m^2l^2-(m\cdot l)^2 .
\]
因此所求的二次点为
\[
f^{ij}=l^2g^{ij}-l^il^j,\qquad\text{即}\qquad f=l^2g-ll,
\]
其中 $l^i=g^{ij}l_j$。验证：对任意线矢 $m$
\[
f^{ij}m_im_j=l^2m^2-(l\cdot m)^2=(m\times l)^2,
\]
并且
\[
f^{ij}l_j=l^2l^i-l^i(l\cdot l)=0,
\]
说明 $f$ 确实退化，$l$ 在其核中。这个张量完全由 $l$ 与 $g$ 确定，而与两个交点分开无关，称为
\textbf{联立二次点}。
\end{proof}

\begin{remark}
$f_{ij}=l^2g_{ij}-l_il_j$ 给出的是 $A,B$ 处两条切线构成的退化二次点列，有时它比 $f^{ij}$ 更
便于应用。此外，$f$ 可以看作 $g$ 在直线 $l$ 上的诱导度规。
\end{remark}

下面用解析坐标看这一等式的意义。取 $a^i=(x_1,y_1,1)$，$b^i=(x_2,y_2,1)$，则
\[
(ab)=\begin{pmatrix}
2x_1x_2 & x_1y_2+x_2y_1 & x_1+x_2\\
x_1y_2+x_2y_1 & 2y_1y_2 & y_1+y_2\\
x_1+x_2 & y_1+y_2 & 2
\end{pmatrix}.
\]
矩阵的元素全部由二次方程的两根之和与两根之积表示。虽然二次方程的解需要开根号，但两根之和与两根
之积不需要，这正是韦达定理。因此等式 $f=l^2g-ll$ 是韦达定理的一个推广。

\begin{remark}[对角元的系数]
写出分量矩阵时容易漏掉对角元上的系数。按对称积的定义 $(ab)^{ij}=a^ib^j+a^jb^i$，当 $i=j$
时对角元为 $a^ib^i+a^ib^i=2a^ib^i$，应带系数 $2$。上表中的 $2x_1x_2,\,2y_1y_2,\,2$ 正是
按此写出的。
\end{remark}

为避免出现不便应用的交叉项 $x_1y_2+x_2y_1$，引入 $k_1=\frac{y_1}{x_1}$，
$k_2=\frac{y_2}{x_2}$，即两点与原点连线的斜率。上式可以改写成
\[
(ab)=\begin{pmatrix}
2x_1x_2 & x_1x_2(k_1+k_2) & x_1+x_2\\
x_1x_2(k_1+k_2) & 2x_1x_2k_1k_2 & y_1+y_2\\
x_1+x_2 & y_1+y_2 & 2
\end{pmatrix}.
\]
设 $AB$ 的中点为 $M$，$x_M=\frac{x_1+x_2}{2}$，$y_M=\frac{y_1+y_2}{2}$。把最后一元归一（即整体
除以 $2$），得
\[
(ab)=\begin{pmatrix}
x_1x_2 & \frac{x_1x_2}{2}(k_1+k_2) & x_M\\[3pt]
\frac{x_1x_2}{2}(k_1+k_2) & x_1x_2k_1k_2 & y_M\\[3pt]
x_M & y_M & 1
\end{pmatrix}.
\]
于是中点坐标以及 $k_1+k_2$、$k_1k_2$ 都可以从分量中读出。

解析几何中常用的齐次化联立，其目的正是把联立方程写成关于斜率的二次方程，从而得到 $k_1+k_2$ 与
$k_1k_2$。由上面的分量式可见，只要算出联立二次点的分量，这一类问题都可以一并解决。

\section{例题：动弦垂直的包络}

\begin{example}
椭圆 $\frac{x^2}{a^2}+\frac{y^2}{b^2}=1$，$O$ 为原点。一条动直线与椭圆交于 $A,B$ 两点，且
$OA\perp OB$。求该动直线的包络。
\end{example}

\begin{proof}[解]
椭圆的张量及其逆为
\[
g_{ij}=\diag\{1/a^2,\,1/b^2,\,-1\},\qquad
g^{ij}=\diag\{a^2,\,b^2,\,-1\}.
\]
设动直线为 $Ax+By+C=0$，即 $l_i=(A,B,C)$。则
\[
l^i=g^{ij}l_j=(a^2A,\;b^2B,\;-C),\qquad
l^2=l_il^i=a^2A^2+b^2B^2-C^2 .
\]
关于 $l^2$ 的这一表达式可以用来判断直线与椭圆的位置关系：相切时 $l^2=0$，交于两点时 $l^2>0$，
无交点时 $l^2<0$。因此 $l^2=a^2A^2+b^2B^2-C^2$ 称为直线与椭圆联立的判别式，它的地位相当于
关于点与椭圆的判别式 $\frac{x^2}{a^2}+\frac{y^2}{b^2}-1$（后者判断点在椭圆内、椭圆上还是
椭圆外）。

联立二次点为 $f^{ij}=l^2g^{ij}-l^il^j$。直接计算
\[
l^il^j=\begin{pmatrix}
a^4A^2 & a^2b^2AB & -a^2AC\\
a^2b^2AB & b^4B^2 & -b^2BC\\
-a^2AC & -b^2BC & C^2
\end{pmatrix},
\]
\[
l^2g^{ij}=\begin{pmatrix}
a^2(a^2A^2+b^2B^2-C^2) & 0 & 0\\
0 & b^2(a^2A^2+b^2B^2-C^2) & 0\\
0 & 0 & -(a^2A^2+b^2B^2-C^2)
\end{pmatrix},
\]
于是
\[
f^{ij}=\begin{pmatrix}
-a^2(C^2-b^2B^2) & -a^2b^2AB & a^2AC\\
-a^2b^2AB & -b^2(C^2-a^2A^2) & b^2BC\\
a^2AC & b^2BC & -(a^2A^2+b^2B^2)
\end{pmatrix}.
\]
二次点 $f$ 描述两个点 $A,B$。过原点的直线 $m_i=(A',B',0)$ 经过 $A$ 或 $B$ 的条件是
$f^{ij}m_im_j=0$，即
\[
f^{11}A'^2+2f^{12}A'B'+f^{22}B'^2=0 .
\]
该二次方程的两个根为 $A'/B'=-k_1,-k_2$，其中 $k_1,k_2$ 是 $OA,OB$ 的斜率，故两根之积
\[
\frac{f^{22}}{f^{11}}=k_1k_2 .
\]
由 $OA\perp OB$ 得 $k_1k_2=-1$，于是
\[
\frac{b^2(C^2-a^2A^2)}{a^2(C^2-b^2B^2)}=-1,
\]
即
\[
b^2(C^2-a^2A^2)+a^2(C^2-b^2B^2)=0,
\]
化简得
\[
A^2+B^2=\Bigl(\frac{1}{a^2}+\frac{1}{b^2}\Bigr)C^2 .
\]
原点到动直线的距离为
\[
d=\frac{|C|}{\sqrt{A^2+B^2}}
=\frac{|C|}{\dfrac{\sqrt{a^2+b^2}}{ab}\,|C|}
=\frac{ab}{\sqrt{a^2+b^2}},
\]
它与直线的位置无关，是常数。因此动直线恒与以原点为心、半径 $ab/\sqrt{a^2+b^2}$ 的圆相切，
所求包络线就是这个圆。
\end{proof}

\begin{remark}
本例中两处符号容易被弄错，值得记录。其一，$l^i=g^{ij}l_j$ 的第三个分量是
$l^3=g^{3j}l_j=-C$，而非 $C$，因为 $g^{33}=-1$。其二，联立二次点应统一取
$f^{ij}=l^2g^{ij}-l^il^j$；若写成 $l^il^j-l^2g^{ij}$ 则与本节的定义相差一个整体符号，
且此时分量矩阵的 $(3,3)$ 元应为 $-(a^2A^2+b^2B^2)$ 而非 $0$。整体符号不影响
$f^{22}/f^{11}$ 这类比值，但本节一律采用 $f=l^2g-ll$，以保证与极线、对合的记号一致。
\end{remark}

\section{对偶：过外点的两条切线}

与直线和椭圆联立对偶，考虑过椭圆外一点 $A$ 作椭圆 $g$ 的两条切线。这两条切线构成的二次线可以
写成
\[
h_{ij}=a^2g_{ij}-a_ia_j,\qquad\text{即}\qquad h=a^2g-aa,
\]
其中 $a_i=g_{ij}a^j$，$a^2=a_ia^i$。事实上 $h_{ij}a^j=a^2a_i-a_i(a_ja^j)=0$，故 $A$ 在 $h$ 的核
方向中，$h$ 上的两条直线都经过 $A$；由对偶命题，这两条直线正是从 $A$ 引出的两条切线。

\begin{remark}
$h^{ij}=a^2g^{ij}-a^ia^j$ 给出的是两个切点构成的退化二次线束，有时它比 $h_{ij}$ 更便于应用。
\end{remark}

下面写出联立二次线的分量式。设两条切线的方程为 $y=k_1x+b_1$ 与 $y=k_2x+b_2$，对应的线矢为
\[
m_i=(k_1,-1,b_1),\qquad n_i=(k_2,-1,b_2),
\]
则 $h_{ij}=m_in_j+m_jn_i$，即
\[
h_{ij}=\begin{pmatrix}
2k_1k_2 & -(k_1+k_2) & k_1b_2+k_2b_1\\
-(k_1+k_2) & 2 & -(b_1+b_2)\\
k_1b_2+k_2b_1 & -(b_1+b_2) & 2b_1b_2
\end{pmatrix}.
\]
一般只会用到前四个分量，使用时注意归一化。

\section{例题：蒙日圆}

\begin{example}
椭圆 $\frac{x^2}{a^2}+\frac{y^2}{b^2}=1$，$C$ 为椭圆外一点。过 $C$ 作椭圆的两条切线互相垂直，
求 $C$ 的轨迹。
\end{example}

\begin{proof}[解]
设 $c^i=(x,y,1)$，则
\[
c_i=g_{ij}c^j=\Bigl(\frac{x}{a^2},\;\frac{y}{b^2},\;-1\Bigr),\qquad
c^2=c_ic^i=\frac{x^2}{a^2}+\frac{y^2}{b^2}-1 .
\]
所求二次线为 $h_{ij}=c^2g_{ij}-c_ic_j$。其中
\[
c_ic_j=\begin{pmatrix}
\dfrac{x^2}{a^4} & \dfrac{xy}{a^2b^2} & \dfrac{x}{a^2}\\[8pt]
\dfrac{xy}{a^2b^2} & \dfrac{y^2}{b^4} & \dfrac{y}{b^2}\\[8pt]
\dfrac{x}{a^2} & \dfrac{y}{b^2} & 1
\end{pmatrix},
\]
逐项相减得
\[
h_{ij}=\begin{pmatrix}
\dfrac{1}{a^2}\Bigl(\dfrac{y^2}{b^2}-1\Bigr) & -\dfrac{xy}{a^2b^2} & -\dfrac{x}{a^2}\\[10pt]
-\dfrac{xy}{a^2b^2} & \dfrac{1}{b^2}\Bigl(\dfrac{x^2}{a^2}-1\Bigr) & -\dfrac{y}{b^2}\\[10pt]
-\dfrac{x}{a^2} & -\dfrac{y}{b^2} & -\Bigl(\dfrac{x^2}{a^2}+\dfrac{y^2}{b^2}\Bigr)
\end{pmatrix}.
\]
两条切线都经过 $C$，以 $C$ 为原点、坐标轴与原来平行时，两条切线的方程之积的二次项部分为
\[
k_1k_2X^2-(k_1+k_2)XY+Y^2,
\]
它与 $h_{ij}$ 的二次项部分成比例，因此斜率之积 $k_1k_2$ 等于 $h_{11}/h_{22}$。由两切线垂直得
$k_1k_2=-1$，即 $h_{11}+h_{22}=0$：
\[
\frac{1}{a^2}\Bigl(\frac{y^2}{b^2}-1\Bigr)+\frac{1}{b^2}\Bigl(\frac{x^2}{a^2}-1\Bigr)=0 .
\]
化简得
\[
x^2+y^2=a^2+b^2 .
\]
因此 $C$ 的轨迹是以原点为心、半径 $\sqrt{a^2+b^2}$ 的圆，称为\textbf{蒙日圆}。
\end{proof}

\begin{remark}
本节的两道例题互为对偶：把点与线的地位互换，一道例题就变成另一道。圆与椭圆在射影几何层面并
没有区别，例中由椭圆恰好得到一个圆，原因在于圆可以用垂直结构来定义。蒙日圆不依赖坐标系的
表达式将在讨论垂直结构时给出。
\end{remark}
